\begin{claim}{D.2.28}{Proved}{The once-suspended functional has only one possible quadratic coefficient}
Retain $B=K(\FF_2,2)$, $A=Q\Sigma B$, $Y=\Sigma^2B$,
$Z=\operatorname{fib}(Y\to K(\FF_2,4))$, and the single
normalization of the relative cup class in \textup{D.2.27}.
Put $W=\Omega Z$ and
$\vartheta_7=\sigma\vartheta_8\in H^7W$.
For every ordinary pointed map
\[
 g:\Sigma B\longrightarrow W
\]
one has $g^*\vartheta_7=0$.

For an arbitrary map $h:\Sigma A\to Y$, choose a lift to
$Z$ and let $G:A\to W$ be its adjunct. Let
$\alpha_3\in H^3A$ be the fundamental class and
$\beta_4=\operatorname{Sq}^1\alpha_3$.
Then, with the same fixed normalization for all maps,
\[
 G^*\vartheta_7=\epsilon\alpha_3\beta_4,
 \qquad
 \theta_h=\sigma(\epsilon\alpha_3\beta_4)
       \quad\text{for some }\epsilon\in\FF_2.
 \tag{D.2.28a}
\]
The possible value $\epsilon=1$ is not excluded here.
Thus this is a complete reduction to one coefficient,
not a vanishing theorem for every once-suspended attachment
or an assertion about subsequent nullification stages.
\end{claim}
\begin{proof}
All coefficients are $\FF_2$. Let
$\lambda_9=\sigma\lambda_{10}\in H^9W$.
The square term in \textup{(D.2.27b)} vanishes under
cohomology suspension, giving the exact identity
\[
 \operatorname{Sq}^4\vartheta_7+
       \operatorname{Sq}^2\lambda_9=0.
 \tag{D.2.28b}
\]
There is a second relation in lower degrees.
Let $\mathcal R_8\in H^8W$ be the unique class whose
cohomology suspension is dual to
$R=\{a,c\}\in H_7F$, where $F=\Omega W$.
The low path-fibration ranges in \textup{D.2.27} make
$H^8W\to H^7F$ an isomorphism and
$H^9W\to H^8F$ injective.
The full degree-eight fiber table gives
\[
 \operatorname{Sq}^2e_6^*=J_1^*+J_3^*,
 \qquad \operatorname{Sq}^1R^*=J_3^*.
 \tag{D.2.28c}
\]
For the second identity, $J_3$ is the only basis class
whose Bockstein has $R$ coefficient one.
For the first, $J_1,J_3$ are exactly the basis classes
whose second square has $e_6$ coefficient one.
Suspension compatibility and the displayed injectivity
therefore prove the actual relation
\[
 \lambda_9=\operatorname{Sq}^2\vartheta_7+
               \operatorname{Sq}^1\mathcal R_8.
 \tag{D.2.28d}
\]
Substitute this into \textup{(D.2.28b)} and use the Adem
identity $\operatorname{Sq}^2\operatorname{Sq}^2
=\operatorname{Sq}^3\operatorname{Sq}^1$:
\[
 \bigl(\operatorname{Sq}^4+
       \operatorname{Sq}^3\operatorname{Sq}^1\bigr)
       \vartheta_7
       +\operatorname{Sq}^2\operatorname{Sq}^1
          \mathcal R_8=0\quad\text{in }H^{11}W.
 \tag{D.2.28e}
\]
This proof uses only the established low suspension
interfaces and the already delooped relation. It does not
replace $H^{11}W$ by an indecomposable quotient or omit
possible higher bar-filtration classes.

Let $u,v,w$ be the Serre generators of degrees $2,3,5$,
with $v=\operatorname{Sq}^1u$ and
$w=\operatorname{Sq}^2v$.
For a map $g:\Sigma B\to W$ write, in the complete
cohomology bases of the source,
\[
 \begin{split}
 g^*\vartheta_7&=a\sigma(u^3)+b\sigma(v^2),\\
 g^*\mathcal R_8&=c\sigma(u^2v)+d\sigma(uw).
 \end{split}
 \tag{D.2.28f}
\]
Cartan and the polynomial description
\cite[\S2, nos.~8--9]{I_Serre1953} give
\[
\begin{aligned}
 (\operatorname{Sq}^4+\operatorname{Sq}^3
                  \operatorname{Sq}^1)(u^3)&=u^5,\\
 (\operatorname{Sq}^4+\operatorname{Sq}^3
                  \operatorname{Sq}^1)(v^2)&=w^2,\\
 \operatorname{Sq}^2\operatorname{Sq}^1(u^2v)&=0,\\
 \operatorname{Sq}^2\operatorname{Sq}^1(uw)&=w^2+u^2v^2.
\end{aligned}
 \tag{D.2.28g}
\]
In the first row,
$\operatorname{Sq}^4u^3=u^5+u^2v^2$ and
$\operatorname{Sq}^3\operatorname{Sq}^1u^3
=\operatorname{Sq}^3(u^2v)=u^2v^2$.
The last row follows from
$\operatorname{Sq}^1(uw)=vw+uv^2$.
Pulling \textup{(D.2.28e)} back gives
\[
 a\sigma(u^5)+(b+d)\sigma(w^2)
       +d\sigma(u^2v^2)=0.
 \tag{D.2.28h}
\]
These three monomials are independent, so $a=b=d=0$.
This proves the universal unit assertion without first
requiring the map to extend to $A$.

For the full source, the complete basis of $H_7A$ is
\[
 Q^4A_0,\quad A_0B_0,\quad F_7,\quad X_7,
 \tag{D.2.28i}
\]
where $A_0=\sigma a_2$, $B_0=\sigma b_3$,
$F_7=\sigma((u^3)^\vee)$, and
$X_7=\sigma((v^2)^\vee)$.
The unit assertion applied to $G|_{\Sigma B}$ kills
the last two coefficients of $G^*\vartheta_7$.
Set $C_5=\sigma((u^2)^\vee)$.
The source Nishida identities give
\[
 \operatorname{Sq}_*^2Q^6C_5=0,
 \qquad
 \operatorname{Sq}_*^4Q^6C_5=Q^4A_0.
 \tag{D.2.28j}
\]
For the second square the leading coefficient is
$\binom42=0$ modulo two and $\beta C_5=0$.
For the fourth square only the term
$Q^4\operatorname{Sq}_*^2C_5=Q^4A_0$ survives.
These are operations inside the free infinite-loop source,
using \cite[Chapters I and III]{I_CohenLadaMay1976};
no ordinary map is assumed to preserve Dyer--Lashof operations.
Pairing \textup{(D.2.28b)} after pullback with $Q^6C_5$
kills the first coefficient in \textup{(D.2.28i)}.

The cohomology class dual to the remaining monomial $A_0B_0$
is exactly $\alpha_3\beta_4$. The unit classes and the
Dyer--Lashof generator in \textup{(D.2.28i)} are primitive,
whereas the reduced diagonal of $A_0B_0$ is
$A_0\otimes B_0+B_0\otimes A_0$.
This proves the first equality in \textup{(D.2.28a)}.
Its second equality is naturality of cohomology suspension
for the actual lift to $Z$ and the universal class of
\textup{D.2.27}.

The retained coefficient cannot be discarded by these
relations alone. In the source write
$\gamma_5=\operatorname{Sq}^2\alpha_3$.
The pair
\[
 q_7=\alpha_3\beta_4,
       \qquad r_8=\alpha_3\gamma_5
 \tag{D.2.28k}
\]
satisfies \textup{(D.2.28e)} algebraically: positive cup
squares vanish in $H^*A$, so
$\operatorname{Sq}^1q_7=0$, and Cartan gives
$\operatorname{Sq}^2\operatorname{Sq}^1r_8
=\operatorname{Sq}^4q_7$.
This is only compatibility of cohomology classes, not
realization by a map. It explains precisely why the
proof stops at the single unresolved coefficient.
\end{proof}
